% !TeX root = ../../main.tex

% Synthesize, benefit, name

In the previous chapters, we discussed two instances of feedback loops over distributed BSP
state.
CARP implements an online but fixed-function loop, reconstructing the global key distribution to
drive adaptive, in-situ range partitioning.
The AMR study implements an offline loop, where trace analysis informs reconfiguration and reruns.
In both cases, the relevant state is not knowable a priori, and making progress requires
iterating between measurement and action.
The AMR experience highlights the practical limitation of doing so offline: iteration latency
becomes the bottleneck, and the resulting insights and mitigations are often deployment-specific
rather than readily transferable.
These properties motivate generalizing the pattern into a substrate that can materialize feedback
via in-situ reductions of application state and support online interventions in response.
Such a capability directly addresses the visibility requirements identified in
\S\ref{sec:amr:concl}: always-on, in-production visibility that surfaces high-level symptoms of
pathologies and enables online steering.
It also subsumes CARP's renegotiation loop, which becomes a policy running on the substrate
rather than a bespoke mechanism.
We call the resulting paradigm \emph{steerable observability}.


Steerable observability requires two capabilities: feedback, materializing views over distributed
state, and control, coordinating interventions across ranks.
The loop may be steered by human operators, automated agents, or policies.
No existing tool provides these as an always-on, BSP-aware capability.
Metrics and logs provide baseline visibility but rarely enough workload structure to attribute
causes.
BSP traces are emitted as per-rank logs and analyzed post-hoc.
Distributed tracing as evolved through
X-Trace~\cite{Fonse2007:X-Trace} and Dapper~\cite{Sigel2010:Dapper} captures causality in
asynchronous microservices by propagating request IDs, but the model does not fit BSP.
Timesteps separated by global synchronization, not requests, are the primary unit of work, and
even identifying the straggler at a single barrier requires aggregating telemetry from all ranks
participating in that barrier.
Control poses a coordination challenge: interventions must be applied with well-defined
consistency and atomicity across ranks.
Parallel debuggers~\cite{:TotalView} provide interactive control but are
designed for debugging sessions, not always-on, in-production observability.


Steerable observability requires a programmable interface for streaming reductions over
distributed telemetry.
Dataframe-style trace analytics are increasingly the norm in both HPC and ML
training workflows~\cite{:HTA,Yildi2023:IOMax,Boehm2016:Caliper,Devar2024:DFTracer}
A streaming, in-situ variant is the logical online counterpart.
We propose \emph{timestep dataframes} as the core abstraction for BSP telemetry: a logical
representation comprising all telemetry emitted during a single timestep across all ranks.
Because BSP synchronization causally partitions execution, each timestep dataframe is a natural
unit of analysis: it summarizes what happened at workload scale and provides the local context
needed to attribute causes.
External stream processing systems can theoretically window event
streams~\cite{Kreps2011:Kafka,Carbo2015:ApacheFlink}, but they decompose this structure into
individual events that must be reaggregated, cannot leverage RDMA fabrics, and cannot provide
coordinated control semantics.
An on-fabric tree-based overlay over columnar data is the natural architecture: Arrow's contiguous
memory layout enables both efficient in-situ reductions and efficient transfers over RDMA, while
the same tree topology supports broadcast for coordinated control.
The result is observability infrastructure that leverages a supercomputer's resources
rather than requiring a second one to process its output.

We present a \emph{tree-based overlay network} (TBON) for real-time observability and control of
large-scale BSP applications, called
ORCA~(\emph{Observability~with~Real-time~Control~and~Aggregation}).
ORCA treats telemetry as columnar streams, transformed via a SQL-based analytics interface called
\emph{OrcaFlow} en route to \emph{sinks} such as dashboards or parallel filesystems.
An RDMA-based TBON-optimized transport (\S\ref{sec:orca:des_xport}) enables low-overhead collection
of timestep dataframes.
They can then be transformed via SQL-based analytics, routed to a dashboard, or persisted as
partitioned Parquet for zero-ETL post-mortems.
This accelerates observability workflows across the spectrum, ranging from real-time monitoring
over application-specific metrics, to offline trace analytics.
Distributed query planning with operator pushdown minimizes data movement across the overlay
(\S\ref{sec:orca:des_insitu}).
A control plane (\S\ref{sec:orca:des_ctl}) provides atomic and timestep-consistent control
semantics, guaranteeing correctness for interventions such as new analytics, probes, and parameter
updates.
ORCA combines a modern columnar analytics
stack~\cite{:ApacheArrow,:ApacheParquet,Lamb2024:DataFusion} with RDMA and BSP-aligned data and
control paths to realize low-overhead, real-time, closed-loop feedback.
This accelerates observability workflows across the spectrum, from application-level metrics
monitored in real time, to post-hoc trace analyses backed by in-situ transforms and
timestep-partitioned Parquet.
% This accelerates observability workflows across the spectrum, ranging from application-level
% metrics monitored in real-time, to post hoc analyses via in-situ filters and transforms and
% partitioned Parquet.

% with collected telemetry that filter or transform telemetry on the fabric and analyze the
% remainder offline, and conventional trace collection persisted as timestep-partitioned Parquet for
% zero-ETL post-mortem analysis.

We evaluate ORCA on a real AMR code at up to 4096 CPU ranks.
% , and progressively study its transport, analytics, and control capabilities.

\begin{itemize}[leftmargin=*]
  \item In conventional tracing workflows, ORCA incurs \textasciitilde{2}\% runtime overhead
        (vs 56--81\% for state-of-the-art tracers), produces up to 42$\times$ smaller
        traces for equivalent telemetry, and enables up to
        6900$\times$ faster analyses than existing trace formats (\S\ref{sec:orca:eval-trace}).
  \item In-situ analytics enable low-overhead needle-in-the-haystack
        workflows---outlier \mpiw{} calls are isolated in real-time at sub-2\% runtime
        overhead, while discarding up to 99.999\% of data at MPI ranks (\S\ref{sec:orca:eval-insitu}).
  \item In an agentic evaluation of its closed-loop capabilities using an LLM
        agent as a proxy for a human operator, ORCA enables localizing a performance
        anomaly---the same bug from our AMR study that took months of manual
        investigation---in 27 minutes, while reducing data volume by up to
        144$\times$ (\S\ref{sec:orca:eval-agent}).
\end{itemize}
More broadly, ORCA demonstrates that declarative SQL analytics can be embedded on a supercomputer
fabric alongside a live MPI application, enabling streaming analytics workflows beyond
observability.
ORCA is open source~\cite{Jain2026:orca-code}.

% This requires a programmable interface for streaming reductions, which must also be declarative to
% support on the fly stuff. Offline trace analysis using dataframe-based analytics. Streaming in
% situ SQL-based over Arrow is a natural online counterpart. We propose an abstraction called
% timestep dataframe for BSP telemetry: similar to Spark RDDs, logical constructs consisting of all
% telemetry emitted from a single timestep. BSP timesteps are causally partitioned: a timestep
% dataframe is a naturally independent unit of analysis as it contains both the global state: what
% happened, and the local state: why it happened. While streaming event processing pipelines can
% theoretically assemble these windows, they will 1. break these nice timestep dataframes into
% separate rows which need to be reassembled afterwards, paying unnecessary cost, 2. cannot leverage
% capabilities like RDMA, 3. can not provide control semantics. As a result, handling
% supercomputer-scale telemetry requires supercomputer-scale analytics clusters incurring
% unnecessary data movement costs. An on-fabric tree-based overlay is a natural choice: enabling
% both hierarchical reductions and broadcast and operator pushdown. Arrow complements RDMA very
% well. Can pushdown. Can use accelerators on path (future capability). Reduction topologies and
% tree-based overlays are common in HPC in the past. On-fabric columnar native TBON is the right
% choice.


% Both CARP and AMR feature materialized views of the application state as inputs into better
% decision-making processes. The process was automatic in CARP but fixed-function. The AMR version
% was manual. Materialized application views periodically occur as a primitive because they are
% targeted renders. They compress a cluster with trillions of operations into specific questions
% asked by a policy or an operator. CARP and AMR studies show that fast views enable fast loops and
% slow views perpetuate bottlenecks. Generalizing these into a programmable substrate allows
% AMR-style studies to become a commonplace operational capacity instead of specialized
% undertakings. Beyond feedback, this has other benefits: by targeting views, we are not
% accumulating raw data all the time. We collect legible signals and use them to intervene on
% demand. No useless data management problem.

% We argue for always on steerable observability. This capability largely does not exist. Metrics
% are nice and captured but also limited. Workload-level insight is missing. Dapper is a useful
% analog. Equivalent semantics were never defined for BSP. ML work is building ad hoc pipelines for
% custom requirements. A common substrate is possible. We show that timesteps and not requests are
% the causally independent units. Telemetry must aggregated by timesteps in a streaming manner via
% programmable interface. Closing feedback loops also requires control semantics that must be
% likewise respect this strucure.

% While event-based streaming systems exist (ref: Kafka and Flink) and can technically aggregate
% streams, they cannot provide control, and they do not understand BSP environment. The execution
% environment is a tightly coupled supercomputer. Telemetry is produced in synchronous bursts.
% Converting it into a signal requires a massive reduction. Even the simplest problem of knowing
% which rank is the straggler at a timestep requires joins over telemetry from all ranks. If the
% problem is externalized by shipping all telemetry to an external system, you have created a
% supercomputer class problem. The smarter thing instead is to adapt to the environment. Use the
% application resourecs and the fabric. Streaming in-situ analytics and control. Our previosu work
% also shows columnar analytics as a recurring feature. Arrow, pushdown and RDMa over a tree based
% overlay are the right solution shape.

% Now we present ORCA. P4 (copy verbatim)
